\chapter{Bites and Stings}

\section*{Definition and Overview}
Bites and stings are injuries caused by animals, insects, and other creatures. They range from trivial mosquito bites and bee stings, which resolve on their own, to serious snake, spider, or marine envenomations and animal bites that carry a high risk of infection.

In Australia, the most important concerns are venomous snakes and spiders, jellyfish and other marine creatures, and infection from dog, cat, and human bites. Prompt first aid and knowing when to seek urgent medical care are the keys to managing most bites and stings safely.

\section*{Epidemiology}
Insect bites and stings are extremely common, and the great majority do not require medical attention.

\begin{itemize}
    \item \textbf{Animal bites:} dog bites are the most frequent animal bites treated in emergency departments, and children are often the victims
    \item \textbf{Bee and wasp stings:} cause allergic reactions in some people; bee stings account for more deaths in Australia than any other venomous creature, usually through anaphylaxis
    \item \textbf{Snake bites:} uncommon but potentially serious, with most occurring in warmer months and in rural or bushland areas
    \item \textbf{Allergy:} around 1 in 50 people may have a significant allergic reaction to insect venom
    \item \textbf{Marine stings:} common among swimmers and surfers, particularly during summer
\end{itemize}

\section*{Aetiology and Pathophysiology}
The effects of a bite or sting depend on the creature involved and the amount of venom or tissue damage.

\begin{itemize}
    \item \textbf{Mechanical injury:} teeth and claws tear the skin and deeper tissues, creating a wound and a route for infection
    \item \textbf{Insect venom:} bee, wasp, and ant venoms cause local pain, swelling, and redness; in allergic people they can trigger anaphylaxis
    \item \textbf{Spiders:} redback spider venom acts on nerves, causing pain, sweating, and sometimes systemic effects, while funnel-web spider venom can be life-threatening in severe cases
    \item \textbf{Snake venom:} Australian elapid venoms affect blood clotting, muscles, and nerves and can cause paralysis, bleeding, and kidney damage
    \item \textbf{Infection:} animal mouths contain many bacteria, and cat bites in particular can inject bacteria deep into the tissues
    \item \textbf{Allergy:} some people develop antibodies to insect venom, so a later sting can trigger a severe reaction
\end{itemize}

\section*{Clinical Features}
\begin{itemize}
    \item \textbf{Local effects:} pain, redness, swelling, itching, and a small puncture wound or raised welt
    \item \textbf{Bites from mammals:} deep puncture wounds, torn skin, bruising, and tenderness; later swelling or discharge may signal infection
    \item \textbf{Snake bite:} often surprisingly little pain at the site; possible bleeding from the gums or wound, muscle weakness, drooping eyelids, vomiting, and difficulty breathing
    \item \textbf{Funnel-web spider bite:} severe pain, sweating, profuse salivation, muscle twitching, and confusion
    \item \textbf{Redback spider bite:} intense local pain that spreads along the limb, with sweating and sometimes nausea
    \item \textbf{Anaphylaxis:} widespread hives, swelling of the lips or throat, wheeze, dizziness, and collapse, which constitutes a medical emergency
    \item \textbf{Delayed signs of infection:} fever, increasing pain, and spreading redness developing days after the injury
\end{itemize}

\section*{Differential Diagnosis}
\begin{itemize}
    \item Simple insect bite versus allergic reaction versus anaphylaxis
    \item An animal bite with infection versus one without
    \item Cellulitis, abscess, or lymphangitis arising from an infected wound
    \item Snake bite versus spider bite versus an unknown envenomation
    \item Reactions to plants, chemicals, or skin conditions such as dermatitis or hives
    \item Tetanus, which must be considered in any dirty or deep wound
    \item Rare conditions such as necrotising fasciitis arising from an infected bite
\end{itemize}

\section*{When to Seek Medical Care}
See a doctor for any bite that is deep, dirty, or from a cat, a human, or an animal of unknown vaccination status; for bites that do not heal; or for signs of infection such as increasing pain, redness, swelling, or fever. Seek medical assessment for any snake bite, funnel-web spider bite, or marine envenomation even if the person feels well, because serious effects can be delayed.

\textbf{Call 000 (or local emergency services) for:}
\begin{itemize}
    \item Any suspected snake bite, especially with vomiting, blurred vision, bleeding, or weakness
    \item A funnel-web spider bite, or any spider bite with severe symptoms
    \item Signs of anaphylaxis: difficulty breathing, throat tightness, dizziness, or collapse
    \item A bite from a large marine animal, or a sting with severe pain, breathing difficulty, or collapse
\end{itemize}

\section*{Diagnosis and Investigations}
\begin{itemize}
    \item \textbf{History:} what bit or stung, when and where it happened, and any past allergic reactions
    \item \textbf{Examination:} assessment of the wound, identification of fang marks, tracking of swelling and redness, and a check of vital signs
    \item \textbf{Blood tests:} for suspected snake bite, including clotting studies and blood counts, with repeated tests as needed
    \item \textbf{Urine tests:} detection of muscle breakdown products after some envenomations
    \item \textbf{Wound swabs:} for suspected infection
    \item \textbf{Allergy assessment:} blood or skin testing after a significant insect venom reaction
    \item \textbf{Imaging:} X-ray or ultrasound if a foreign body or deep tissue involvement is suspected
\end{itemize}

\section*{Management}
\begin{itemize}
    \item \textbf{First aid for snake and funnel-web bites:} keep the person calm and still, immobilise the affected limb, and apply a pressure-immobilisation bandage from the bite site upward
    \item \textbf{For most other bites and stings:} wash the area, apply a cold pack, and take simple pain relief
    \item \textbf{Antivenom:} given in hospital for significant snake, spider, and some marine envenomations
    \item \textbf{Bee stings:} scrape the stinger out with a fingernail or flat object; do not squeeze it
    \item \textbf{Tick removal:} freeze or carefully remove the whole tick without squeezing its body
    \item \textbf{Allergic reactions:} antihistamines for mild reactions; adrenaline via auto-injector for anaphylaxis, followed by emergency care
    \item \textbf{Wound care:} cleaning, debridement, and antibiotics for infected or high-risk bites
    \item \textbf{Tetanus and rabies considerations:} vaccination updated as needed
    \item \textbf{Referral:} to a clinical immunologist for venom immunotherapy after a serious sting reaction
\end{itemize}

\section*{Complications}
\begin{itemize}
    \item Wound infection, including cellulitis, abscess, and occasionally infection of the bone or joints
    \item Cat and human bites carry a particular risk of deep infection and scarring
    \item Anaphylaxis from insect stings, which can be fatal without prompt treatment
    \item Complications of snake bite: bleeding, muscle damage, kidney failure, and nerve injury
    \item Permanent tissue damage or amputation in severe marine envenomations
    \item Tetanus in unvaccinated people
    \item Scarring and psychological fear of the creature involved
\end{itemize}

\section*{Prognosis and Outcomes}
Most bites and stings heal quickly without treatment or with simple first aid. With prompt hospital care and antivenom, the outlook for snake and spider envenomation is good, and deaths are rare in Australia.

Allergic people who carry and use an adrenaline auto-injector markedly reduce their risk from future stings. Wound infections usually settle with antibiotics, although cat and human bites may require surgical cleaning. Recovery from a serious envenomation can take weeks and may leave muscle or kidney damage in some cases.

\section*{Prevention and Self-Care}
\begin{itemize}
    \item Wear shoes and long trousers when walking in bushland, and avoid reaching into rock crevices or long grass
    \item Shake out boots, clothing, and bedding that have been left outdoors in snake and spider country
    \item Do not try to catch or kill snakes; back away slowly
    \item Keep food and rubbish sealed to avoid attracting animals
    \item Avoid provoking dogs and other animals, and supervise children near them
    \item If allergic to stings, carry an adrenaline auto-injector and a written action plan, and wear a medical alert identifier
    \item Avoid perfumes and bright clothing in bee territory, and do not disturb hives
    \item Keep tetanus vaccinations up to date
\end{itemize}

\section*{Sources and Further Reading}
The following organisations provide reliable information on bites and stings for patients, families, and health professionals.

\begin{itemize}
    \item NHS. \textit{Information on insect bites and stings and their treatment.} https://www.nhs.uk/
    \item Mayo Clinic. \textit{Overview of insect bites and stings, and when to seek care.} https://www.mayoclinic.org/
    \item World Health Organization. \textit{Resources on snakebite envenoming and its prevention.} https://www.who.int/
    \item MedlinePlus. \textit{Health information on bites and stings for patients.} https://medlineplus.gov/
    \item Centers for Disease Control and Prevention. \textit{Resources on animal and insect bites.} https://www.cdc.gov/
\end{itemize}
